%% ============================================================
%% Introduction to Cognitive Psychology — Week 9 Study Notes
%% ============================================================
\documentclass[11pt,a4paper]{article}

\usepackage[a4paper, left=1.8cm, right=1.8cm, top=1.3cm, bottom=1.8cm]{geometry}
\usepackage{booktabs}
\usepackage{tabularx}
\usepackage{array}
\usepackage{enumitem}
\usepackage{titlesec}
\usepackage{fancyhdr}
\usepackage{fontenc}
\usepackage{inputenc}
\usepackage{microtype}
\usepackage{parskip}
\usepackage{hyperref}
\usepackage{amsmath}
\usepackage{amssymb}

\hypersetup{hidelinks, colorlinks=false}

%% ── Table helpers ─────────────────────────────────────────
\newcolumntype{L}[1]{>{\raggedright\arraybackslash}p{#1}}

%% ── Section headings ──────────────────────────────────────
\titleformat{\section}{\large\bfseries}{}{0em}{}[\titlerule]
\titleformat{\subsection}{\normalsize\bfseries}{}{0em}{}
\titleformat{\subsubsection}{\normalsize\bfseries}{}{0em}{}

\titlespacing{\section}{0pt}{12pt}{4pt}
\titlespacing{\subsection}{0pt}{8pt}{3pt}
\titlespacing{\subsubsection}{0pt}{6pt}{2pt}

%% ── Page style ────────────────────────────────────────────
\pagestyle{fancy}
\fancyhf{}
\renewcommand{\headrulewidth}{0pt}
\fancyfoot[C]{\small Cognitive Psychology --- Week 9 \textbar\ Page \thepage}

%% ── Misc helpers ──────────────────────────────────────────
\newcommand{\bb}[1]{\textbf{#1}}
\setlength{\parskip}{4pt}

%% ═══════════════════════════════════════════════════════════
\begin{document}
\setlength{\arrayrulewidth}{0.4pt}

\begin{center}
\bfseries\LARGE Introduction to Cognitive Psychology\par
\vspace{0.2cm}
\large Assignment 9 Detailed Solutions
\end{center}
\vspace{0.5cm}

%% ════════════════════════════════════════════════════════════
\section*{ASSIGNMENT 9 --- Full Solutions with Explanations}

\subsubsection*{Question 1}
\textit{Which of the following is/are generally agreed upon as necessary criteria/criterion for claiming that a communication system is a language?}
\medskip

\begin{tabular}{L{0.3cm} L{14.9cm}}
& A.\ Regularity \\
& B.\ Productivity \\
& C.\ Referentiality \\
& \textbf{D.\ Both regularity and productivity \checkmark} \\
\end{tabular}

\vspace{0.3cm}
\noindent\textbf{Ans:} D — Both regularity and productivity\\[0.2cm]
\textbf{Why this answer:} \\
Both \textbf{regularity} and \textbf{productivity} are generally agreed upon as necessary criteria for a communication system to qualify as a \textbf{language}. \textbf{Regularity} means that language follows systematic rules (grammar, syntax, phonology) that govern how symbols are combined. \textbf{Productivity} (also called generativity) means that speakers can produce and understand an \textbf{unlimited number of novel sentences} using a finite set of rules and symbols. While referentiality (the ability to refer to things) is also a feature of language, regularity and productivity together are the most commonly cited defining criteria.
\vspace{0.4cm}

\subsubsection*{Question 2}
\textit{Infinite combinations of ideas can be expressed in language. In other words, language is:}
\medskip

\begin{tabular}{L{0.3cm} L{14.9cm}}
& A.\ Regular \\
& B.\ Interpersonal \\
& \textbf{C.\ Productive \checkmark} \\
& D.\ Referential \\
\end{tabular}

\vspace{0.3cm}
\noindent\textbf{Ans:} C — Productive\\[0.2cm]
\textbf{Why this answer:} \\
\textbf{Productivity} (or generativity) is the property of language that enables speakers to produce and comprehend an \textbf{infinite number of ideas and sentences} using a \textbf{finite set of symbols and rules}. This is one of the hallmarks that distinguishes human language from other animal communication systems. Any speaker can create entirely new sentences that have never been uttered before and be understood by other speakers of the same language.
\vspace{0.4cm}

\subsubsection*{Question 3}
\textit{A fundamental problem of speech perception, according to Miller, is that:}
\medskip

\begin{tabular}{L{0.3cm} L{14.9cm}}
& A.\ Speech is discrete rather than continuous \\
& \textbf{B.\ A single phoneme sounds different depending on its context \checkmark} \\
& C.\ Hearing is a less accurate rather than vision \\
& D.\ Missing phonemes can render words incomprehensible \\
\end{tabular}

\vspace{0.3cm}
\noindent\textbf{Ans:} B — A single phoneme sounds different depending on its context\\[0.2cm]
\textbf{Why this answer:} \\
According to Miller, a fundamental challenge in speech perception is the phenomenon known as \textbf{coarticulation} — the fact that the acoustic realization of a \textbf{single phoneme varies significantly} depending on the surrounding phonemes (its phonetic context). For example, the ``d'' sound in ``dog'' is acoustically different from the ``d'' sound in ``bed.'' This context-dependent variability, known as the \textbf{lack of invariance problem}, makes it difficult to explain how listeners reliably identify phonemes despite their constantly changing acoustic properties.
\vspace{0.4cm}

\subsubsection*{Question 4}
\textit{The branch of linguistics devoted to the study of meaning:}
\medskip

\begin{tabular}{L{0.3cm} L{14.9cm}}
& A.\ Phonetics \\
& \textbf{B.\ Semantics \checkmark} \\
& C.\ Morphology \\
& D.\ Pragmatics \\
\end{tabular}

\vspace{0.3cm}
\noindent\textbf{Ans:} B — Semantics\\[0.2cm]
\textbf{Why this answer:} \\
\textbf{Semantics} is the branch of linguistics specifically devoted to the study of \textbf{meaning} — both the meaning of individual words (lexical semantics) and the meaning of sentences and larger units of discourse (compositional semantics). It is distinct from \textbf{phonetics} (the study of speech sounds), \textbf{morphology} (the study of word structure and formation), and \textbf{pragmatics} (the study of how context influences the interpretation of meaning in communication).
\vspace{0.4cm}

\subsubsection*{Question 5}
\textit{Did she say ``many'' or ``men knee''? This type of ambiguity is referred to as:}
\medskip

\begin{tabular}{L{0.3cm} L{14.9cm}}
& \textbf{A.\ Phonetic \checkmark} \\
& B.\ Lexical \\
& C.\ Syntactic \\
& D.\ Semantic \\
\end{tabular}

\vspace{0.3cm}
\noindent\textbf{Ans:} A — Phonetic\\[0.2cm]
\textbf{Why this answer:} \\
\textbf{Phonetic ambiguity} occurs when the same continuous stream of speech sounds can be \textbf{segmented or interpreted in different ways}, leading to different words or phrases. The example ``many'' vs.\ ``men knee'' illustrates this perfectly — the same acoustic signal can be parsed as either a single word or two separate words. This is related to the \textbf{segmentation problem} in speech perception, where listeners must determine where one word ends and another begins in the continuous flow of speech.
\vspace{0.4cm}


\subsubsection*{Question 6}
\textit{The underlying knowledge that allows people to produce and comprehend their language is called:}
\medskip

\begin{tabular}{L{0.3cm} L{14.9cm}}
& A.\ Linguistic performance \\
& \textbf{B.\ Linguistic competence \checkmark} \\
& C.\ Linguistic production \\
& D.\ Grammatical behaviour \\
\end{tabular}

\vspace{0.3cm}
\noindent\textbf{Ans:} B — Linguistic competence\\[0.2cm]
\textbf{Why this answer:} \\
\textbf{Linguistic competence}, a concept introduced by Noam Chomsky, refers to the \textbf{underlying, internalized knowledge} of a language that a speaker possesses — including knowledge of grammar, vocabulary, and the rules for combining them. This is distinguished from \textbf{linguistic performance}, which refers to the actual \textbf{use of language} in real situations, which may include errors, hesitations, and other imperfections. Competence represents what you \textit{know} about your language; performance represents how you \textit{use} it.
\vspace{0.4cm}
\newpage
\subsubsection*{Question 7}
\textit{The study of the ways in which sounds can be combined in any given language is called:}
\medskip

\begin{tabular}{L{0.3cm} L{14.9cm}}
& \textbf{A.\ Phonology \checkmark} \\
& B.\ Morphology \\
& C.\ Syntax \\
& D.\ Grammar \\
\end{tabular}

\vspace{0.3cm}
\noindent\textbf{Ans:} A — Phonology\\[0.2cm]
\textbf{Why this answer:} \\
\textbf{Phonology} is the branch of linguistics that studies the \textbf{rules governing how sounds (phonemes) are organized and combined} in a particular language. It examines which sound combinations are permissible and which are not. For example, in English, ``bl'' is a permissible initial consonant cluster (as in ``block''), but ``bn'' is not. Phonology is distinct from \textbf{morphology} (study of word formation), \textbf{syntax} (study of sentence structure), and \textbf{grammar} (a broader term encompassing multiple levels of linguistic rules).
\vspace{0.4cm}

\subsubsection*{Question 8}
\textit{We notice ambiguities in sentences:}
\medskip

\begin{tabular}{L{0.3cm} L{14.9cm}}
& A.\ All the time \\
& \textbf{B.\ In ``garden path'' sentences \checkmark} \\
& C.\ When they are humorous \\
& D.\ When they make no sense at all \\
\end{tabular}

\vspace{0.3cm}
\noindent\textbf{Ans:} B — In ``garden path'' sentences\\[0.2cm]
\textbf{Why this answer:} \\
\textbf{Garden path sentences} are sentences that initially lead the reader or listener toward one syntactic interpretation, only to force a \textbf{reanalysis} when later information contradicts that initial parse. A classic example is: ``The horse raced past the barn fell.'' These sentences make us \textbf{acutely aware of ambiguity} because our initial interpretation fails, and we must consciously re-parse the sentence. In everyday language processing, we typically resolve ambiguities automatically without noticing them; it is only in garden path sentences that the ambiguity becomes salient.
\vspace{0.4cm}

\subsubsection*{Question 9}
\textit{According to research, which of the following sets of sentences would take the longest amount of time to read and comprehend?}
\medskip

\begin{tabular}{L{0.3cm} L{14.9cm}}
& A.\ We took the dog to the vet. The dog was nervous. \\
& B.\ Sara drove the car to the market. The car ran out of gas. \\
& \textbf{C.\ She wrapped the Christmas presents. The sweater needed a larger box. \checkmark} \\
& D.\ We got some beer out of the car. The beer was warm. \\
\end{tabular}

\vspace{0.3cm}
\noindent\textbf{Ans:} C — She wrapped the Christmas presents. The sweater needed a larger box.\\[0.2cm]
\textbf{Why this answer:} \\
This question relates to \textbf{bridging inferences} in language comprehension. In options A, B, and D, the second sentence directly refers back to the same entity mentioned in the first sentence (the dog, the car, the beer), making the connection \textbf{immediately explicit}. In option C, however, the reader must make an \textbf{additional inference}: that ``the sweater'' was one of the ``Christmas presents'' being wrapped. This requires an extra cognitive step — a bridging inference — which takes \textbf{more processing time}. Research shows that sentences requiring such inferences are read more slowly because the reader must construct the implicit connection between the two sentences.
\vspace{0.4cm}
\newpage

\subsubsection*{Question 10}
\textit{Speech acts in which the utterance itself is the action — such as ``You're fired!'' — are considered to be which type of speech act?}
\medskip

\begin{tabular}{L{0.3cm} L{14.9cm}}
& \textbf{A.\ Declaration \checkmark} \\
& B.\ Directive \\
& C.\ Commissive \\
& D.\ Expressive \\
\end{tabular}

\vspace{0.3cm}
\noindent\textbf{Ans:} A — Declaration\\[0.2cm]
\textbf{Why this answer:} \\
A \textbf{declaration} (also called a \textbf{declarative speech act}) is a type of speech act in which the \textbf{utterance itself brings about a change in reality}. When someone with the appropriate authority says ``You're fired!'' the act of saying it \textit{is} the act of firing — the words themselves constitute the action. Other examples include ``I now pronounce you husband and wife'' and ``I hereby resign.'' This is distinct from \textbf{directives} (requests or commands), \textbf{commissives} (promises or commitments), and \textbf{expressives} (expressions of feelings or attitudes).
\vspace{0.4cm}

\medskip
\subsection*{Assignment Quick Answer Summary}

\begin{center}
\renewcommand{\arraystretch}{1.4}
\begin{tabular}{|L{0.8cm}|L{5.5cm}|L{8.9cm}|}
\hline
\bb{Q\#} & \bb{Topic} & \bb{Correct Answer} \\
\hline
1 & Language Criteria & D — Both regularity and productivity \\
2 & Language Productivity & C — Productive \\
3 & Speech Perception Problem & B — Phoneme sounds differ by context \\
4 & Linguistics \& Meaning & B — Semantics \\
5 & Phonetic Ambiguity & A — Phonetic \\
6 & Linguistic Competence & B — Linguistic competence \\
7 & Sound Combination Study & A — Phonology \\
8 & Sentence Ambiguity & B — In ``garden path'' sentences \\
9 & Longest Reading Time & C — Christmas presents / sweater \\
10 & Declarative Speech Acts & A — Declaration \\
\hline
\end{tabular}
\end{center}

\medskip
\end{document}
